\documentclass[../../main.tex]{subfiles}
\begin{document}

\begin{center}
  {\fontsize{48pt}{58pt}\selectfont \textbf{1989}\normalsize\textbf{年度}}
\end{center}

\begin{table}[htb]
  \centering
  \caption{出題テーマと難易度}
  \begin{tabular}{|c|p{9cm}|c|}
    \hline
    \textbf{問題番号} & \textbf{テーマ} & \textbf{難易度} \\
    \hline
    第1問 & 二等辺三角形の反復構成と極限 & \\
    \hline
    第2問 & 相加平均の条件を満たす数列 & \\
    \hline
    第3問 & 三次式と導関数の除法 & \\
    \hline
    第4問 & 四面体の切断面と面積 & \\
    \hline
    第5問 & 立方体の面の塗り分け確率 & \\
    \hline
    第6問 & 円錐形タンクの注水時間 & \\
    \hline
  \end{tabular}
\end{table}

\begin{table}[htb]
  \centering
  \caption{解答}
  \renewcommand{\arraystretch}{1.4}
  \begin{tabular}{|c|c|p{9cm}|}
    \hline
    \textbf{問題番号} & \textbf{小問} & \textbf{答え} \\
    \hline
    第1問 & (1) & $\dfrac{\sin\theta}{4}$ \\
    \cline{2-3}
    & (2) & $\dfrac{\sin\theta}{\theta}$ \\
    \hline
    第2問 & -- & $a_1=a_2=\cdots=a_n$ \\
    \hline
    第3問 & -- & 実数解はただ1つ（証明問題） \\
    \hline
    第4問 & -- &  \\
    \hline
    第5問 & -- & $P(N)=\dfrac{(N-1)(N-2)}{N^5}$．$a<b$なら$P(a)>P(b)$ \\
    \hline
    第6問 & -- & あと$16$時間 \\
    \hline
  \end{tabular}
\end{table}

\paragraph{試験形態}
6問．試験時間は要確認．

\end{document}
